% ID: SEL-D0024-OB-E01
% ESTADO: validada. S025, cotejada con examen y criterios.
\textbf{a)} La última condición equivale a $x+y\ge4$.
Como $1.1+2.8=3.9<4$, el punto queda fuera.

\textbf{b)} El recinto está por encima de $y=4-x$ y por debajo
de $y=2x+1$ y $y=13-4x$. Al igualar las fronteras dos a dos,
obtenemos los vértices $(1,3)$, $(2,5)$ y $(3,1)$.
\begin{center}
\setlength{\unitlength}{7mm}
\begin{picture}(8,8)(-1,-1)
\qbezier(0,0)(2,0)(4,0)\put(4.2,0){$x$}
\qbezier(0,0)(0,3)(0,6)\put(0,6.2){$y$}
\qbezier(1,3)(1.5,4)(2,5)
\qbezier(2,5)(2.5,3)(3,1)
\qbezier(3,1)(2,2)(1,3)
\put(1,3.2){\makebox(0,0)[r]{$(1,3)$}}
\put(2,5.2){$(2,5)$}
\put(3.2,1){$(3,1)$}
\end{picture}
\end{center}
En esos vértices $F$ vale $1.5$, $1.5$ y $-7.5$, respectivamente.
El mínimo $\boxed{-7.5}$ se alcanza sólo en $(3,1)$.
Como $F=1.5(y-2x)$, el máximo $\boxed{1.5}$ se alcanza en todo
el lado $y=2x+1$ con $1\le x\le2$.

\textbf{c)} $F=0$ equivale a $y=2x$. Por ejemplo, $(2,4)$
cumple $4\le5$, $4\le5$ y $2+4\ge4$, por lo que pertenece al recinto.
